%%%%%%%%%%%%%%%%%%%%%%%%%%%%%%%%%%%%%%%%%
% Beamer Presentation
% LaTeX Template
% Version 1.0 (10/11/12)
%
% This template has been downloaded from:
% http://www.LaTeXTemplates.com
%
% License:
% CC BY-NC-SA 3.0 (http://creativecommons.org/licenses/by-nc-sa/3.0/)
%
%%%%%%%%%%%%%%%%%%%%%%%%%%%%%%%%%%%%%%%%%

%----------------------------------------------------------------------------------------
%	PACKAGES AND THEMES
%----------------------------------------------------------------------------------------

\documentclass{beamer}

\mode<presentation> {

% The Beamer class comes with a number of default slide themes
% which change the colors and layouts of slides. Below this is a list
% of all the themes, uncomment each in turn to see what they look like.

%\usetheme{default}
%\usetheme{AnnArbor}
%\usetheme{Antibes}
%\usetheme{Bergen}
%\usetheme{Berkeley}
%\usetheme{Berlin}
%\usetheme{Boadilla}
%\usetheme{CambridgeUS}
%\usetheme{Copenhagen}
%\usetheme{Darmstadt}
%\usetheme{Dresden}
%\usetheme{Frankfurt}
%\usetheme{Goettingen}
%\usetheme{Hannover}
\usetheme{Ilmenau}
%\usetheme{JuanLesPins}
%\usetheme{Luebeck}
%\usetheme{Madrid}
%\usetheme{Malmoe}
%\usetheme{Marburg}
%\usetheme{Montpellier}
%\usetheme{PaloAlto}
%\usetheme{Pittsburgh}
%\usetheme{Rochester}
%\usetheme{Singapore}
%\usetheme{Szeged}
%\usetheme{Warsaw}

% As well as themes, the Beamer class has a number of color themes
% for any slide theme. Uncomment each of these in turn to see how it
% changes the colors of your current slide theme.

%\usecolortheme{albatross}
%\usecolortheme{beaver}
%\usecolortheme{beetle}
%\usecolortheme{crane}
%\usecolortheme{dolphin}
%\usecolortheme{dove}
%\usecolortheme{fly}
%\usecolortheme{lily}
%\usecolortheme{orchid}
%\usecolortheme{rose}
%\usecolortheme{seagull}
%\usecolortheme{seahorse}
%\usecolortheme{whale}
%\usecolortheme{wolverine}

%\setbeamertemplate{footline} % To remove the footer line in all slides uncomment this line
%\setbeamertemplate{footline}[page number] % To replace the footer line in all slides with a simple slide count uncomment this line

%\setbeamertemplate{navigation symbols}{} % To remove the navigation symbols from the bottom of all slides uncomment this line
}

\usepackage{graphicx} % Allows including images
\usepackage{booktabs} % Allows the use of \toprule, \midrule and \bottomrule in tables

%----------------------------------------------------------------------------------------
%	TITLE PAGE
%----------------------------------------------------------------------------------------

\title[Summary Presentation : Paper 16]{Extrapolating human judgments from skip-gram vector representations of word meaning (Hollis et al., 2017)} % The short title appears at the bottom of every slide, the full title is only on the title page

\author{Shawon Ashraf} % Your name
\institute[IMS] % Your institution as it will appear on the bottom of every slide, may be shorthand to save space
{
Universität Stuttgart \\ % Your institution for the title page
%\medskip
%\textit{john@smith.com} % Your email address
}
\date{25 January, 2022} % Date, can be changed to a custom date

\begin{document}

\begin{frame}
\titlepage % Print the title page as the first slide
\end{frame}

%\begin{frame}
%\frametitle{Overview} % Table of contents slide, comment this block out to remove it
%\tableofcontents % Throughout your presentation, if you choose to use \section{} and \subsection{} commands, these will automatically be printed on this slide as an overview of your presentation
%\end{frame}

%----------------------------------------------------------------------------------------
%	PRESENTATION SLIDES
%----------------------------------------------------------------------------------------

%------------------------------------------------
\section{Goals and Motivations} % Sections can be created in order to organize your presentation into discrete blocks, all sections and subsections are automatically printed in the table of contents as an overview of the talk
%------------------------------------------------

\begin{frame}
\frametitle{Goals of the paper}
\begin{itemize}
    \item Predict / Extrapolate Human Judgement from vector representation of words
        \begin{itemize}
            \item Using lexical and semantic properties
        \end{itemize}
    \item How different are ratings from statistical models and actual human judgement?
        \begin{itemize}
            \item Especially, ratings for affective variables: valence, arousal, dominance
        \end{itemize}
    \item Are large norm sets useful for Psycholinguistic research?
    \item Provide new norms with larger coverage of text
\end{itemize}
\end{frame}

%------------------------------------------------

\begin{frame}
\frametitle{Motivation}
\begin{itemize}
\item Existing statistical models can not capture lexical and semantic properties
\item Predicting ratings from large amount of data is tedious and time consuming
    \begin{itemize}
        \item Need automated methods
        \item Combine large scale data and statistical models to predict lexical judgement
    \end{itemize}
\end{itemize}
\end{frame}

%------------------------------------------------
\section{Methods and Findings}
\begin{frame}
\frametitle{Methods}
\begin{block}{Setup}
    \begin{itemize}
        \item Skipgram vectors from Word2Vec
        \item Linear Regression 
        \item Feature pruning with backward regressing
        \item Forward regressing with AIC criterion
        \item Use norms from ANEW dataset
        \item Evaluate predicted ratings against human ratings
        \item Sentiment based classification for affective variables
    \end{itemize}
\end{block}

\end{frame}

%------------------------------------------------

\begin{frame}
\frametitle{Findings}
\begin{itemize}
    \item Statistical models can judge lexical properties similar to humans
    \item However, it is difficult for affective variables
    \item Large norm sets can help finding new research areas in Psycholinguistics
\end{itemize}
\end{frame}

%------------------------------------------------
\section{End}
%------------------------------------------------

%------------------------------------------------

\begin{frame}
\Huge{\centerline{Thank you!}}
\end{frame}

%----------------------------------------------------------------------------------------

\end{document} 